\section{Conclusion}
\label{sec:conclusion}

\method{} reframes voice anti-cloning as a purification-aware protection problem. 
Instead of assuming that a protected waveform is fed directly into a cloning system, we model the stronger and increasingly realistic attacker who first removes perturbations with a purification front end. 
For diffusion purifiers, protection must survive not only direct speaker extraction but also the attacker's reverse denoising path; 
for traditional signal-processing front ends, it must also survive lightweight transformations such as resampling, quantization, filtering, and denoising.

\method{} addresses this gap by combining pre-purification speaker disruption, post-purification speaker disruption, and trajectory-level deviation inside the purifier, 
and by providing a benchmark protocol with clean-purification controls and calibrated evaluator thresholds. 
The planned ablations are designed to test whether trajectory supervision is a necessary complement to pre- and post-purification identity optimization, and whether output-level gradients alone fail under stochastic purification.

The broader implication is that voice protection should be evaluated as part of an adaptive purification-and-cloning chain, not only as a direct perturbation against a cloning backend. 
We plan to provide an anonymized implementation and evaluation protocol to support reproducible review of this problem.
